\cotefile{59}
\scanpage{15}
\mspage{14}

e) Généralisation des constructions précédentes.

On revient aux conditions de a), mais on suppose donné un complexe \deleted{parfait} $F^{\bullet}$ sur $X$, \deleted{\unsure{et un complexe}} parfait (ou seulement parfait rel $/S$) et on considère

\[
\det{}^{*}\mathbb R\operatorname{pr}_{2*}(\operatorname{pr}_1^{*}(F^{\bullet})\otimes\mathscr L_g)=M_g
\]

On aura, pour un autre choix $g'$, \deleted{et} un isom. can.

\[
M_{g'}=M_g\,g'_{\unsure{P|p}}{}^{*}(L_g)^{-\chi_F}=M_g(N_gN_{g'}^{-1})^{\chi_F}
\]

où maintenant $\chi$ est la fonction rang pour $\mathbb R\operatorname{pr}_{2*}(\operatorname{pr}_1^{*}(F)\otimes\mathscr L_g)$. On trouve alors, pour $F$ fixé, un homomorphisme

\[
\varphi_{F^{\bullet}}\colon\underline{\mathrm{NS}}_{X/S}\longrightarrow\underline{\mathrm{NS}}_{B/S}
\]

dépendant additivement de $F^{\bullet}$ (pour la structure de catégorie \deleted{\unsure{additive}}\inserted{triangulée} de $D_{\mathrm{parf}}(X)$ resp $D_{\mathrm{parf}}(X/S)$).

On trouve alors, pour tte section $\delta$ de $\underline{\mathrm{NS}}_{X/S}$, un hom

\[
X\xrightarrow{{\unsure{l|h|\ell}}_{\delta}^{F^{\bullet}}}\underline{\mathrm{Pic}}^{\varphi_F(\delta)}_{B/S}\overset{A}{\times}\bigl({\unsure{P|p}}^{\delta}\overset{B}{\times}(A,\varphi_F\widetilde{(\delta)})\bigr)
\]

d'où un hom de $A$-torseurs

\[
\mathrm{Alb}^{-\chi_F(\delta)}_{X/S}\mathbin{\unsure{-|\to}}\underline{\mathrm{Pic}}^{\varphi_F(\delta)}_{B/S}\times\bigl(\underline{\mathrm{Pic}}^{\delta}_{X/S}\overset{B}{\times}(A,\varphi_F\widetilde{(\delta)})\bigr)
\]

et enfin une section

\[
\boxed{L_F(\delta)\in\Gamma\bigl(S,\underline{\mathrm{Alb}}_{X/S}^{\chi_F(\delta)}\overset{A}{\times}\underline{\mathrm{Pic}}_{B/S}^{\varphi_F(\delta)}\times\bigl(\underline{\mathrm{Pic}}_{X/S}^{\delta}\overset{B}{\times}(A,\varphi_F\widetilde{(\delta)})\bigr)\bigr)}
\]

Il faut enfin calculer $\chi_F(\delta)$ et $\varphi_F(\delta)$, et pour ceci, comme en \unsure{c)|e)}, on est ramené au cas $S=\operatorname{Spec} k$, $k$ corps alg. clos, et à calculer

\[
\chi(L_\delta\otimes F^{\bullet})\quad\text{et}\quad\det{}^{*}\mathbb R\operatorname{pr}_{2*}\bigl(\mathscr L\otimes\operatorname{pr}_1^{*}(L_\delta\otimes F)\bigr)
\]

On note que dans les calcul, $\delta$ \unsure{disparait|disparaît}, à condition de remplacer $F$ par $\deleted{F\otimes}L_\delta\otimes F$ \ldots\unsure{.}
\scanpage{16}
\mspage{15}

On peut le faire par RR si $X$ est loc. une intersection complète (rel $/S$), on trouve

\[
\begin{aligned} &\deleted{\unsure{\chi(L_\delta\otimes F)|\chi(L_\delta)}}\\ \chi_F(\delta)&=\chi(L_\delta\otimes F^{\bullet})=\operatorname{tr}\bigl(\operatorname{ch}(L_\delta\,\deleted{\unsure{\otimes F}}{\unsure{\otimes|\,}}F)\operatorname{Todd}(X)\bigr)\\ &\deleted{\unsure{\varphi}}\qquad=\text{terme constant dans }\underbrace{\operatorname{pr}_{2*}\bigl(\operatorname{pr}_1^{*}(\operatorname{ch}_X(L_\delta\otimes F)\inserted{\operatorname{Todd}(X)}\deleted{\unsure{\operatorname{Todd}}})\,\operatorname{ch}_{X\times B}\unsure{\exp(D)|\exp D}\bigr)}_{(\exp D)\,.[\operatorname{ch}_X(L_\delta\otimes F)\operatorname{Todd}(X)]}\end{aligned}
\]

\[
\varphi_F(\delta)=\text{terme de degré 1 dans }\;(\exp D)\,.[\operatorname{ch}_X(L_\delta\otimes F)\operatorname{Todd}(X)]
\]

\uline{Ex. 1}\qquad $X$ lisse $/S$ de dim relative 1. On aura

\[
\operatorname{ch}_X(\underset{\delta\,\text{de degré }d}{L_\delta}\otimes F)=(1+\underset{c(L_\delta)}{\delta})(\underset{\text{rang } F^{\bullet}}{\unsure{\rho|p}}+\underset{\det^{\prime}(F)}{\gamma})=1+\bigl(\underset{\substack{\text{degré }d\\\text{classe }\delta}}{\delta}+\underset{\substack{\text{degré }c\\\deleted{\text{\unsure{catégorie} }C}}}{\gamma}\bigr)
\]

\[
\operatorname{Todd}(X)=1+\underset{\text{degré }1-g}{\underline{\tfrac12 K}}
\]

\[
\operatorname{ch}_X(L_\delta\otimes F)\otimes\operatorname{Todd}(X)=1+\tau,\qquad\text{avec}\qquad\deg\tau=1-g+d+c
\]

On \unsure{a|--} \deleted{Il faut \unsure{prendre le terme|prendre les termes}}

\[
\boxed{\chi_F(\delta)=1-g+d+c}
\]

D'autre part, $\varphi(\delta)$ n'a pas changé \deleted{\unsure{pour}}\inserted{ici} par rapport au calcul précédent dans \unsure{c)|e)}

\[
\boxed{\varphi_F(\delta)=\varphi(\delta)=\tfrac12\operatorname{pr}_{2*}(D^2)=\text{polarisation canonique de }B}
\]

Donc on trouve une section canonique de

\[
\underline{\mathrm{Alb}}_{X/S}^{\chi_F(\delta)=\chi(\delta)+c}\overset{A}{\times}\mathrm{Pic}_{B/S}^{\varphi(\delta)}\overset{A}{\times}\bigl(\mathrm{Pic}_{X/S}^{\deleted{\unsure{\varphi(}}\delta\deleted{\unsure{)}}}\overset{B}{\times}(B,\widetilde{\varphi(\delta)})\bigr)
\]

donc, en utilisant le résultat fait dans c, on trouve une section canonique de $\underline{\mathrm{Alb}}^{c}_{X/S}\simeq\underline{\mathrm{Pic}}^{c}_{X/S}$.

Il n'y a pas de mérite à cela, car il suffit de prendre la section définie par \inserted{le Module inversible} $\det{}^*F$, sur $X/S$.
\scanpage{17}
\mspage{16}

\uline{Ex} 2\qquad $X$ un schéma abélien $A$. On a ici $\operatorname{Todd}(X)=1$.

On trouve

\[
\chi_F(\delta)=\operatorname{tr}(\operatorname{ch}(L_\delta F))=\deg a_n
\]

si on pose

\[
\operatorname{ch}(L_\delta F)=\sum_1^n a_i\qquad a_i\in\operatorname{Chow}(A)
\]

On trouve ensuite, comme avant

\[
\varphi_F(\delta)=\inserted{\text{terme de degré 1 de }}(\exp D)\,.(a_{n-1})=\operatorname{pr}_{2*}(\tfrac12\underline{D^2a_{n-1}}).
\]

f)\qquad Nous allons donner une autre interprétation de l'homomorphisme $l_\delta$ de d), en définissant (en intégrant exclusivement sur $X\times_S B$) un homomorphisme

\[
\boxed{l^{\unsure{*|' }}\colon X\times\underline{\mathrm{Pic}}_{X/B}\longrightarrow\underline{\mathrm{Pic}}_{B/S}}.
\]

\margin{(additif en $\underline{\mathrm{Pic}}_{X/S}$)}

Pour ceci, soit $g$ une section de $X$, et $\lambda$ une section de $\underline{\mathrm{Pic}}_{X/S}$, qui (grâce à $g$) est définie par un module inversible $L\inserted{{}=L(\lambda,g)}$ sur \deleted{\unsure{le}} $X$, \unsure{trivialisée|trivialisé} le long de $g$. On peut alors former sur $X\times_S B$ ($W_g\inserted{{}=W(g)}$ faisceau de Weil sur $X\times_S B\unsure{)|^{g}}$, correspondant à $\unsure{\theta|0|g}$) $\det{}^{*}(\operatorname{pr}_{2*}(\operatorname{pr}_1^{*}(L)\otimes W))$ \deleted{\unsure{classe}} faisceau inv. sur $B$, dont la classe \deleted{\unsure{a}}\inserted{définit} une section $l(g,\lambda)$ de $\underline{\mathrm{Pic}}_{B/S}$. La formation est compatible avec extension de la base, d'où l'hom $l$.

Désignons, pour une section $g$ de $X$, par $l_g$ l'hom $\underline{\mathrm{Pic}}_{X/S}\mathbin{\unsure{-|\to}}\underline{\mathrm{Pic}}_{B/S}$ qu'il définit.
\scanpage{18}
\mspage{17}

On va préciser comment $l_g$ varie avec $g$. Soit $g'$ une autre section. \unsure{Alors|Nous} on a

\[
L(\lambda,g')=L(\lambda,g)\,g'{}^{*}(L(\lambda,g))^{-1}=
\]

\[
\text{\deleted{donc}}\qquad W(g')=W(g)\,g'_B{}^*(W(g))^{-1}
\]

(où $g'_B$ est la section de $X\times_S B/B$ déduite de la section $g'$ de $X/S$). On a donc

\[
\operatorname{pr}_1^{*}(L(\lambda,g'))W(g')=\operatorname{pr}_1^{*}(L(\lambda,g'))W(g)\,D
\]

\margin{$\begin{aligned}D={}&g'_B{}^{*}(L(\lambda,g))^{-1}\otimes_{\unsure{\mathcal O_X|O_X}}\\ &g'_B{}^*(W(g))^{-1}\end{aligned}$}

d'où comme d'habitude

\[
\begin{aligned}l_{g'}(\lambda)&=l_g(\lambda)+\bigl[\text{classe dans }\mathrm{Pic}_{B/S}\text{ de }g'_B{}^{*}(W(g))^{-\chi(\lambda)}\\ &\hspace{12em}=(N_{g'}N_g^{-1})^{-\chi(\lambda)}\\ &=l_g(\lambda)+(\psi_X(g')-\psi_X(g))(-\chi(\lambda))\end{aligned}\tag*{(*)}
\]

où $\psi_X\colon X\longrightarrow\underline{\mathrm{Alb}}^1_{X/S}$ est l'hom canonique, et $\underline{\mathrm{Alb}}^1_{X/S}$ est considéré comme torseur sous $A=\underline{\mathrm{Pic}}^0_{B/S}$ de la façon habituelle. Donc, considérant l'hom

\[
\varphi\colon\underline{\mathrm{NS}}_{X/S}\longrightarrow\underline{\mathrm{NS}}_{B/S}
\]

induit par $l$ de façon évidente, on déduit pour tte section $\delta$ de $\underline{\mathrm{NS}}_{X/S}$ un hom

\[
\boxed{l_\delta\colon X\times\underline{\mathrm{Pic}}^\delta_{X/S}\longrightarrow\underline{\mathrm{Pic}}^{\varphi(\delta)}_{B/S}}
\]

\deleted{qui, \unsure{sur} définit \unsure{; se donne}} \deleted{se factorise $\underline{\mathrm{Alb}}^{\deleted{\unsure{-\chi(\delta)|\chi(\delta)}}}_{X/S}$ d'ailleurs par $\underline{\mathrm{Alb}}^{-\chi(\delta)}_{X/S}\overset{A}{\times}$} satisfaisant \unsure{aux,|à,|les,}
\scanpage{19}
\mspage{18}

De plus, pour une section $g$ de $X/S$, on vérifie que l'on a pour $\beta\deleted{\unsure{\in A}}\in\underline{\mathrm{Pic}}^0_{X/S}(S)=B(S)$

\[
l_g\bigl(\underset{\underset{\underline{\mathrm{Pic}}^\delta_{X/S}(S)}{\rotatebox{90}{$\in$}}}{\lambda}+\beta\bigr)=l_g(\lambda)\,\deleted{\unsure{+}}\unsure{-|+}\,\widetilde{\varphi(\delta)}(\beta)\tag*{**}
\]

\deleted{En effet, on sait que (comme tt \unsure{$S$-hom.|flèche} de} \deleted{torseurs sous des schémas abéliens différents)}

\[
\deleted{l_g^\delta\colon\underset{\text{torseur sous }B}{\underline{\mathrm{Pic}}^\delta_{X/S}}\longrightarrow\underset{\text{torseur sous }A}{\underline{\mathrm{Pic}}^{\varphi(\delta)}_{B/S}}\qquad\text{est associé}}
\]

\deleted{à un hom $\rho\colon B\longrightarrow A$, et il faut prouver que} \deleted{$\rho=-\widetilde{\varphi(\delta)}$. On peut pour ceci supposer $S$ le} \deleted{spectre d'un corps alg. clos\inserted{ $k$}, et il suffit de} \deleted{prouver que l'on a égalité pour les pts à} \deleted{valeurs dans $k$, i.e. \inserted{prouver} la formule ** pour $S=\operatorname{Spec}k$.} \deleted{Il suffit de la prouver pour des $\alpha$ engendrant} \deleted{le groupe $B(k)$, et donc pour $\alpha$ de la forme} \deleted{$\psi(g')-\psi(g)$ $(g',g\in X)$. \unsure{Choisissons un}}

On est ramené, grâce au fait que le faisceau inv. $L_\beta$ sur $X$ provient d'un faisceau \deleted{\unsure{unité associé} $L_\beta$}\inserted{$W_\beta$} sur $A$ (via $X\mathbin{\unsure{-|\to}}A$ déduit de $g$), sur \deleted{faisceau $W_\beta$ sur}\inserted{celui-ci défini par}
% slanted

\margin{\unsure{Vérif.|Réf.} ; Attention aux conventions de \unsure{signe|signes}}
le faisceau de Weil $W$ sur $A\times B$, au cas où $X=A$, $g=$ section nulle. Avec les notations développées ailleurs, on doit donc \unsure{identifier|expliciter|vérifier}

\[
\det\underline{\unsure{\Phi|\mathscr L|\mathfrak F}}_{\mathscr L}(L\otimes L_\beta)\simeq\det\Phi_{\mathscr L}(L)\otimes\underbrace{L'_{-\varphi(\delta)(\beta)}}_{\substack{\text{faisceau sur }B\text{ associé}\\\text{à }-\varphi(\delta)(\beta)}}\,,
\]
\scanpage{20}
\mspage{19}

Or on a

\[
\deleted{\det}\Phi_{\mathscr L}(L\otimes L_\beta)=\Phi_{\mathscr L}(L)*\varepsilon_{-\beta}
\]

d'où

\[
\det\Phi_{\mathscr L}(L\unsure{\otimes|{}}L_\beta)\simeq\deleted{\unsure{\Phi}}\det\Phi_{\mathscr L}(L)*\varepsilon_{-\beta}\deleted{\unsure{=}}
\]

donc la classe de $\det(\Phi_{\mathscr L}(L\otimes L_\beta))$ dans $\underline{\mathrm{Pic}}_{B/S}(S)$ est égale à \deleted{$\det($} celle de $\det\Phi_{\mathscr L}(L)+\widetilde{\varphi(\delta)}(-\beta)$
% slanted

\margin{Vérif, avec Attention au \unsure{signe|signes}\ldots.}
(où $\varphi(\delta)$ \unsure{c'|est} l'image de $\det\Phi_{\mathscr L}(L)$ dans $\mathrm{NS}(B)$), d'où la conclusion annoncée.

Par suite, on peut interpréter $l_g$ comme un hom. $\underline{\mathrm{Pic}}^\delta_{X/S}\overset{B}{\times}(A,-\widetilde{\varphi(\delta)})$ dans $\underline{\mathrm{Pic}}^{\varphi(\delta)}_{B/S}$, ou encore comme une \deleted{\unsure{hom.}} section de $\underline{\mathrm{Pic}}^\delta_{X/S}\overset{B}{\times}(A,\widetilde{\varphi(\delta)})\overset{A}{\times}\underline{\mathrm{Pic}}^{\varphi(\delta)}_{B/S}$. On trouve ainsi \deleted{\unsure{$N$}}\inserted{un} hom (pour $g$ variable)

\[
X\mathbin{\unsure{-|\to}}\underline{\mathrm{Pic}}^{\varphi(\delta)}_{B/S}\overset{A}{\times}\underline{\mathrm{Pic}}^\delta_{X/S}\overset{B}{\times}(A,\unsure{\overline{\varphi(\delta)}|\widetilde{\varphi(\delta)}})
\]

qui n'est autre que \deleted{celui de} l'hom $l_\delta$ de d), comme on vérifie facilement.

\uline{Remarques}\qquad a) On n'a pas eu, dans cette interprétation, à faire intervenir un faisceau sur $\underline{\mathrm{Pic}}^\delta_{X/S}$, -- seult sur $B$.

Notons que si $\underline{\mathrm{Pic}}^0_{X/S}$ n'est pas un schéma abélien, on peut encore donner un sens aux constructions précédentes, en supposant \uline{donné} un schéma abélien
\scanpage{21}
\mspage{20}

$B$ sur $S$ et un hom de groupes

\[
B\xrightarrow{i}\underline{\mathrm{Pic}}_{X/S},
\]

ce qui définit d'ailleurs un torseur $Q$ sous $A=\underline{\mathrm{Pic}}^0_{B/S}$ et un hom canonique

\[
X\xrightarrow{\psi} Q
\]

[C'est par un \unsure{sorite|suite} sur la correspondance entre des $i$ et des $\psi$ que devrait commencer \unsure{mon|un|une} \deleted{\unsure{le}} petit topo de Albanese\ldots.].

Tout ce qui a été dit \unsure{vaudra|vaut} alors, on trouve

\[
X\times\underline{\mathrm{Pic}}_{X/S}\mathbin{\unsure{-|\to}}\underline{\mathrm{Pic}}_{B/S}
\]

définissant

\[
X\times\underline{\mathrm{Pic}}^\delta_{X/S}\longrightarrow\underline{\mathrm{Pic}}^{\varphi(\delta)}_{B/S}
\]

\margin{\uline{NB} : \unsure{sous-entendus|sous-entendre} dans la notation $\varphi(\delta)$}

qui donne, en fait, si on a un $B$\deleted{\unsure{-teur|-torseur}} $P$ et un hom $P\longrightarrow\underline{\mathrm{Pic}}^\delta_{X/S}$ \inserted{compatible avec $i$}

\[
X\longrightarrow\underline{\mathrm{Pic}}^{\varphi(\delta)}_{B/S}\overset{A}{\times}\bigl(\deleted{\underline{\mathrm{Pic}}_{X/S}^{\delta}}P\overset{B}{\times}(A,\widetilde{\varphi(\delta)})\bigr)
\]

\deleted{\unsure{et}}\inserted{se factorisant par}

\[
Q^{-\chi(\delta)}\longrightarrow\underline{\mathrm{Pic}}^{\varphi(\delta)}_{B/S}\overset{A}{\times}\bigl(\deleted{\underline{\mathrm{Pic}}^\delta_{X/S}}P\overset{B}{\times}(A,\widetilde{\varphi(\delta)})\bigr)
\]

i.e. on a une section de

\[
Q^{\chi(\delta)}\overset{A}{\times}\underline{\mathrm{Pic}}^{\varphi(\delta)}_{B/S}\overset{A}{\times}\bigl(\deleted{\underline{\mathrm{Pic}}^\delta_{X/S}}P\overset{B}{\times}(A,\widetilde{\varphi(\delta)})
\]

\deleted{b) On peut aller plus loin, dans la \deleted{\unsure{réduction}} simplification} \deleted{et généralisation des énoncés, lorsqu'on se borne à} \deleted{des $X$ munis de sections \unsure{marquées} $g$, \deleted{Alors} du} \deleted{fait qu'on a $X\longrightarrow A$. \unsure{Ceci noté}, $X$ n'interv}
\scanpage{23}
\mspage{1}
% Paper clip at the upper left.

Soit $X/S$ propre, plat, de présentation finie, $f_*(\underline{\mathcal O}_X)\simeq\underline{\mathcal O}_S$ \unsure{univ.!|univ.} $\underline{\mathrm{Pic}}^0_{X/S}$ un schéma abélien $B$. Soit $A=\underline{\mathrm{Pic}}^0_{B/S}$. On définit alors un \deleted{\unsure{schéma}} $A$-torseur $\underline{\mathrm{Alb}}^1_{X/S}$, et un $S$-morphisme

\[
\psi\colon X\longrightarrow\underline{\mathrm{Alb}}^1_{X/S},
\]

ou ce qui revient au $=$, un homomorphisme

\[
\delta\colon X\times_S X\longrightarrow A\qquad\qquad\delta(x,y)=\psi(y)-\psi(x)
\]

tel que

\[
\delta(x,x)=0\qquad.
\]

Pour ceci, on prend, pour deux sections $x,y$ de $X$ sur $S$

\[
\delta(x,y)=\inserted{\text{classe de }}N_yN_x^{-1}\qquad\text{dans }\underline{\mathrm{Pic}}^0_{B/S}(S)
\]

où $N$ est un module inversible de Weil sur $X\times_S B$ (qui existe grâce à l'existence d'une section de $X$ sur $S$).

On prouve la propriété universelle de $(Q=\underline{\mathrm{Alb}}^1_{X/S},\psi)$ pour les $S$-morphismes de $X$ dans les \deleted{schémas} torseurs sous schémas abéliens (ou plus généralement, sous des schémas en groupes commutatifs quelconques?).

Si on définit $Q'$ par $\delta'=\delta^{-1}$, on trouve un isomorphisme $Q'=Q^{-1}$, $\psi'=\psi^{-1}$, et $(Q',\psi')$ est une autre solution du $=$ Pb universel, canoniquement isomorphe au premier. Donc, lorsqu'on définit $(\underline{\mathrm{Alb}}^1,\psi)$ comme solution du Pb universel, avec $A$ le schéma abélien des translations de $\underline{\mathrm{Alb}}^1$, il y a une \uline{convention de signe} à faire pour choisir un isomorphisme $A\simeq\underline{\mathrm{Pic}}^0_{B/S}$. \uline{Nous prendrons la première convention}
% archive sheet number: 22
\scanpage{24}
\mspage{2}

\uline{de signe}. Supposons alors $X$ \inserted{lisse,} de dim. relative 1. On peut considérer alors $\underline{\mathrm{Pic}}^d_{X/S}$, \inserted{torseur sous $B=\underline{\mathrm{Pic}}^0_{X/S}$,} correspondant aux faisceaux inv sur $X/S$ de degré relatif $d$. Il y a un morphisme canonique

\[
X\xrightarrow{\psi}\underline{\mathrm{Pic}}^1_{X/S},
\]

\deleted{et celui-ci} défini par la diagonale de $X\times_S X$, qui est un diviseur relatif sur le deuxième facteur.

\diagram{\begin{tikzcd}[ampersand replacement=\&,column sep=large,row sep=large] X\arrow[r,no head,"\operatorname{pr}_1"]\arrow[d,no head]\&X\times_S X\arrow[d,no head,"\operatorname{pr}_2"]\\ S\arrow[r,no head]\&X\end{tikzcd}}

Il fait de $\underline{\mathrm{Pic}}^1_{X/S}$ une solution du pb universel d'Albanese. Il s'ensuit, grâce à la première convention de \unsure{signe|signes}, un isom

\[
c\colon B\xrightarrow{\sim} A=\underline{\mathrm{Pic}}^0_{B/S},
\]

qui s'explicite ainsi : Si $y,x$ sont deux sections de $X$ sur $S$, alors $\psi(y)-\psi(x)$ est une section de $\underline{\mathrm{Pic}}^0_{X/S}$ correspondant au diviseur relatif $y(S)-x(S)$, et on a

\[
c(\psi(y)-\psi(x))=c(\operatorname{cl}(y(S))-\operatorname{cl}(x(S))=N_yN_x^{-1}
\]

où $N$ est un faisceau de Weil sur $X\times_S B$.

\uline{Question à tirer au clair} : Cet isom $c$ définit-il \deleted{\unsure{le}} la $\Theta$-\uline{polarisation} sur $B$? \deleted{\unsure{(ou la symétrique d'une}}\inserted{Non, c'est sans doute la} \deleted{\unsure{polarisation)).}}\inserted{\uline{symétrique}} de la $\Theta$-polarisation.

C'est une question ``géométrique,'' donc on peut supposer $S$ spectre d'un corps alg. clos. On a une polarisation de $B$ par \unsure{le|la} \unsure{classique|d'unique|l'unique} diviseur $\Theta$. L'hom \inserted{correspondant} $u\colon B\longrightarrow A$ est défini par $u(\deleted{\unsure{x}}t)=\text{classe de }t.\Theta-\Theta$. Si
% archive sheet number: 23
\scanpage{25}
\mspage{3}

$t$ est de la forme $\operatorname{cl}((y)-(x))$, on aura donc

\[
u((y)-(x))=\text{classe de }\ldots\unsure{.'|.}\;\deleted{\unsure{y.\Theta-x.\Theta-\Theta}}
\]

% archive sheet number: 24
\scanpage{26}
\mspage{4}

\uline{Schémas abéliens polarisés principaux}.

Soit $A/S$ un schéma abélien, muni d'une section $\delta$ de $\underline{\mathrm{NS}}_{A/S}$, d'où un torseur $P^\delta=\underline{\mathrm{Pic}}^\delta_{A/S}$ sous le schéma abélien dual $B=\underline{\mathrm{Pic}}^0_{A/S}$, et un faisceau inversible $\mathscr L$ sur $A\times P^\delta$ \ldots. D'autre part, $\delta$ définit $\widetilde\delta\colon A\longrightarrow B$, et les opérations de $A$ sur $P^\delta$ via $\widetilde\delta$ sont celles déduites par transport de structure des opérations de $A$ sur lui-$=$ par translations.

Supposons que $\delta$ soit une polarisation, \deleted{\unsure{de|on}} alors on sait que les $H^i$ des $\mathscr L_x$ $(x\in P^\delta)$ sont nuls pour $i>0$, et que $\deleted{\unsure{R}}\operatorname{pr}_{2*}(\mathscr L)$ est donc un module loc libre de rang $\chi(\delta)$. Si $\chi(\delta)=1$, c'est un module inversible sur $P^\delta$, et il lui est associé un diviseur canonique $D_\delta$ \deleted{$\unsure{\Theta_\delta|\Theta}$} sur $A\times P^\delta$ relativement à $P^\delta$. D'autre part, $\widetilde\delta\colon A\mathbin{\unsure{-|\to}}B$ est alors un isom, et $P^\delta$ peut être regardé comme un $A$-torseur.

\deleted{\unsure{l'effet d'une translation}} \deleted{\unsure{$\Theta$}} Le diviseur $\deleted{\unsure{\Theta|D}}D_\delta$ est stable par opérations de $A$ opérant diagonalement sur $A\times P^\delta$ \deleted{\unsure{(sur les|Donc il}} Il provient donc d'un diviseur sur le quotient $(A\times P^\delta)/A=P^\delta$, noté \deleted{\unsure{aussi}} $\Theta_\delta$. Ce dernier a la vertu que $\chi(\Theta_\delta)=1$.

\uline{Proposition}. Soit $A$ un schéma abélien $/S$, \inserted{de dim relative $g$}, \inserted{Soit} $\delta$ une polarisation sur $A$ de degré 1. Alors il existe un $A$-torseur
% archive sheet number: 25
\scanpage{27}
\mspage{5}

$P$, et un diviseur \inserted{relatif positif} $\Theta$ sur $P$, [tel que $\deg\Theta^g/g!=1$] et que la polarisation $\delta$ soit égale à celle déduite de la polarisation de $P$ par le diviseur $\Theta$. Etant donné \deleted{\unsure{Et}} $P'$ et $\Theta'$ ayant les \unsure{$=^{s}$|mêmes} vertus que $P$ et $\Theta$, il existe un unique $A$-isom de $P$ sur $P'$, transformant $\Theta$ en $\Theta'$ [en d'autres termes \unsure{toute|seule} \deleted{\unsure{autre}} section de $A$ qui invarie $\Theta$ est la section nulle].

b) Si on a un torseur $P$ sous $A$, muni d'un diviseur relatif positif $\Theta$ \deleted{de degré de polarisation 1,}\inserted{tel que $\deg\Theta^g/g!=1$}, alors $\Theta$ est ample rel $/S$ et définit \deleted{\unsure{donc}} une polarisation de $P/S$ donc \deleted{degré} de $A/S$. Cette dernière est de degré de polarisation 1.

c) Sous les conditions de b), considérons l'hom de composition $A\times P\longrightarrow P$ ; et soit $\Theta'$ l'image inverse de $\Theta$ dans $A\times P$. Alors $\Theta$ est un diviseur positif relatif sur $A\times P\inserted{/P}$, et définit donc un \deleted{\unsure{f}} faisceau inversible sur $A\times P$, donc un hom $P\mathbin{\unsure{-|\to}}\underline{\mathrm{Pic}}_{A/S}$. Ce dernier est un isom de $P$ sur $\underline{\mathrm{Pic}}^\delta_{A/S}$, où $\delta$ est la polarisation de $A/S$ définie par \deleted{$\unsure{\delta}$} $\Theta$ sur $P$.

d) Supposons que $A=\underline{\mathrm{Alb}}^0_{X/S}$, où $X/S$ est une courbe propre et lisse relative de genre $g$. \deleted{Alors} Considérons la polarisation canonique $\delta$ de $A$ \deleted{\unsure{de $B$, d'où un iso}}\inserted{\unsure{définie|défini} par le diviseur diagonal} sur $X\times_S X$. Cette polarisation est de degré 1, et $P^\delta$ est canoniquement isomorphe à \inserted{comme $A$-torseur} $\underline{\mathrm{Pic}}^{g-1}_{X/S}=\underline{\mathrm{Alb}}^{g-1}_{X/S}$, est isomorphe\inserted{ à} \unsure{sous|sur}
% slanted

\margin{question de signe!}
% archive sheet number: 26
\scanpage{28}
\mspage{6}

transformant $\Theta$ en $W_{g-1}$ (ce qui le caractérise).

\begin{center}\rule{1.8em}{0.7pt}\end{center}

\deleted{Soit $X/S$ propre plat de prés. finie, $f_*(\underline{\mathcal O}_X)\simeq\underline{\mathcal O}_S$ \unsure{univ.t|univt|univ.}, $\underline{\mathrm{Pic}}^0_{X/S}$ un} \deleted{schéma abélien $B$, de dual $A=\underline{\mathrm{Alb}}^0_{X/S}$} \deleted{Supposons donné sur \deleted{\unsure{$X\times_S X$}} $(X\times X)$ une correspondance} \deleted{divisorielle \deleted{\unsure{$c$}} $C$ \inserted{symétrique}, définissant donc une correspondance} \deleted{divisorielle \inserted{symétrique} sur $A$ donc une \deleted{polarisation $\gamma$ de $A$}} \deleted{\unsure{section de} $\underline{\mathrm{NS}}_{A/S}$ \unsure{de} soit $\gamma$.}

\deleted{\unsure{e bis}} (\unsure{e|c} bis)\qquad Soit $(P,\Theta)$ comme dans \unsure{b)|B)}. Considérons \deleted{$P^{-1},\Theta^{-1}$} $P^{-1}$, et l'isom (compatible avec la \uline{symétrie} sur $A$) $P\xrightarrow{\sigma_P}P^{-1}$, et soit $\Theta^{\unsure{-1|-*}}=\sigma_P(\Theta)$. Alors la polarisation définie par $\Theta^{-1}$ sur $A$ est égale à celle définie par $\Theta$, donc il y a un isomorphisme unique de $A$-torseurs

\[
K\colon P\xleftarrow{\sim}P^{-1}
\]

tel que $K(\Theta)=\Theta^{-1}$. On peut aussi considérer \unsure{ce|ceci|$K$} comme une section de $P^{\overset A\otimes 2}$

\[
K\in\Gamma(S,P^{\overset A\otimes 2})
\]

% slanted

\margin{Vérifier que c'est la section \unsure{$\delta$} de $(\underline{\mathrm{Pic}}^{\unsure{2\delta|\delta}}_{A/S}$ canonique définie pour toute section $\delta$ de $\underline{\mathrm{NS}}_{A/S}$ \ldots}

Lorsque $A=\underline{\mathrm{Alb}}^0_{X/S}$, $X,\delta$ comme dans d), alors \deleted{$K$ est la} moyennant l'isom. $\inserted{P\simeq}\underline{\mathrm{Pic}}^{g-1}_{X/S}$ \deleted{\unsure{Alb|$\simeq\underline{\mathrm{Alb}}$}}, d'où $P^{(2)}\simeq\underline{\mathrm{Pic}}^{2g-2}_{X/S}$, $K$ est la section provenant du \deleted{le} faisceau $\underline{\Omega}^1_{X/S}$ sur $X$.
% archive sheet number: 27
